\scanpage{11}
on conclut (cf lemme) que $A_{S''}$ est constant. \deleted{Il l'est aussi en tant que variété polarisée, par fonctorialité. Il ne reste le choix que d'un nb fini de rigidifications de Jacobi, \unsure{donc} $(S''_s)_{\mathrm{réd}}\longrightarrow R$ se factorise par $\unsure{\mathrm{Sp}| \operatorname{Spec}}(k(s))$, donc $S''_s$ est réduit à un point.} On gagne.

Le cas où on ne suppose plus $S''$ fini sur $S'$ se traite par \unsure{la méthode|les méthodes} essentiellement de $[\,]$ p 75-76, \uline{à condition} qu'on prouve que l'isomorphisme sur $U''\simeq U$ entre $T_p(A'')|U''$ et $M_{S''}|U''$ se prolonge à $S''$. \deleted{Il faudrait tirer au clair cette question.} \inserted{OK}
\margin{?}

c) Cas $S$ \deleted{normal} \inserted{géom. unibranche}. Choisissons un \unsure{ent|ent.} premier $\ell\neq p,3$, et soit $U'\longrightarrow U$ un morphisme étale fini surjectif \deleted{qui \unsure{rig}} tel que $A_{U'}$ soit muni d'une \unsure{rigigidification|rigidification} d'échelon $\ell$.

% slanted
\margin{\unsure{utilisons|initialisons|inutile} un \unsure{rev.|fond} !}
Normalisons $S$ dans $R(U')/R(U)$ \inserted{$R(S)$}, d'où morphisme fini \inserted{surjectif} $S'\longrightarrow S$, \inserted{prolongeant $U'-U$} avec $S'$ normal intègre. Alors on peut trouver un schéma abélien $A'$ sur $S'$ qui prolonge $(A_U)\times_U U'$, avec $T_p(A')\simeq p^*(M)$.

\diagram{\begin{tikzcd}[ampersand replacement=\&,row sep=small,column sep=small]
 U\arrow[d,arrows={Hooks[right]-}] \& U'\arrow[l]\arrow[d,arrows={Hooks[right]-}] \& U''\arrow[l,shift left]\arrow[l,shift right]\arrow[d,arrows={Hooks[right]-}]\\
 S \& S'\arrow[l,"p"'] \& S''\arrow[l,shift left]\arrow[l,shift right]
\end{tikzcd}}
Alors par la partie ``pl. fid'' appliquée à $S''=S'\times_S S'$, utilisant que $U''$ est dense dans $S''$ (\deleted{\unsure{par Chevalley}} \inserted{N.B.} $S''-S$ est ouvert car $S'-S$ est univ\textsuperscript{t} ouvert \ldots), on trouve

% archive sheet number: 10
